\newpage
\addcontentsline{toc}{chapter}{Appendix 2: Project Code}
\chapter*{Appendix 2: Project Code}
\markboth{Appendix 2: Project Code}{}

\section*{Project Structure}

\begin{verbatim}
Arambh/
|-- frontend/                    # React 18 + Vite Frontend
|   |-- src/
|   |   |-- components/         # Reusable UI components
|   |   |   |-- auth/          # Auth guards (ProtectedRoute, AdminRoute)
|   |   |   |-- chat/          # Chat interface components
|   |   |   |-- layout/        # Navbar, Sidebar, Footer
|   |   |   `-- ui/            # ShadCN/Radix primitives
|   |   |-- hooks/             # Custom React hooks
|   |   |-- layouts/           # PublicLayout, AuthLayout, MainLayout,
|   |   |                        DashboardLayout
|   |   |-- lib/               # API client, i18n config
|   |   |-- pages/             # Route pages (14 pages + admin/superadmin)
|   |   |   |-- admin/        # Admin dashboard, uploads, FAQs, URLs
|   |   |   `-- superadmin/   # SuperAdmin colleges, users, analytics
|   |   |-- store/             # Zustand stores (auth, chat, college,
|   |   |                        theme, sidebar)
|   |   |-- App.jsx            # Main application component
|   |   `-- index.css          # Tailwind CSS styles
|   `-- package.json
|
|-- backend/                     # Python FastAPI Backend Gateway
|   |-- app/
|   |   |-- api/               # API route modules
|   |   |   |-- auth.py       # JWT authentication endpoints
|   |   |   |-- conversations.py # Chat session management
|   |   |   |-- colleges.py   # College CRUD + search
|   |   |   |-- admin.py      # Admin operations
|   |   |   |-- superadmin.py # SuperAdmin management
|   |   |   |-- recommend.py  # College recommendation
|   |   |   `-- health.py     # Health check
|   |   |-- core/             # Config, database, deps, logging
|   |   |-- models/           # Pydantic schemas
|   |   |-- repositories/     # MongoDB data access layer
|   |   |-- services/         # Business logic (auth, AI client)
|   |   |-- scripts/          # Seed scripts
|   |   `-- main.py           # FastAPI application entry point
|   `-- requirements.txt
|
|-- ai-services/                 # LangGraph Multi-Agent RAG Engine
|   |-- app/
|   |   |-- agents/           # LangGraph agent nodes
|   |   |   |-- graph.py     # StateGraph definition
|   |   |   |-- nodes.py     # Agent nodes (query, retrieval,
|   |   |   |                   generation, validation)
|   |   |   |-- runner.py    # Streaming orchestrator
|   |   |   |-- state.py     # AgentState TypedDict
|   |   |   |-- safety.py    # Prompt injection detection
|   |   |   `-- recommend.py # College recommendation agent
|   |   |-- rag/             # RAG pipeline components
|   |   |   |-- retriever.py # Hybrid dense+BM25 retrieval
|   |   |   |-- embeddings.py# HuggingFace embeddings
|   |   |   |-- vectorstore.py # Qdrant operations
|   |   |   |-- bm25_index.py# Sparse BM25 index
|   |   |   |-- reranker.py  # Cross-encoder reranking
|   |   |   |-- ingestion.py # Document ingestion pipeline
|   |   |   |-- chunker.py   # Text chunking
|   |   |   |-- cleaner.py   # HTML/PDF cleaning
|   |   |   |-- web_search.py# Live DuckDuckGo search
|   |   |   |-- translation.py # Language detection + translation
|   |   |   |-- prompts.py   # All LLM prompt templates
|   |   |   `-- llm.py       # Ollama LLM interface
|   |   |-- api/             # FastAPI routes (agent, ingest, health)
|   |   |-- core/            # Config, database, logging
|   |   |-- eval/            # RAGAS evaluation scripts
|   |   `-- main.py          # FastAPI application entry point
|   `-- requirements.txt
|
|-- crawler/                     # Async Web Crawler + Celery Workers
|   |-- app/
|   |   |-- celery_app.py    # Celery configuration
|   |   |-- crawler.py       # Main crawl orchestrator
|   |   |-- fetcher.py       # HTTP fetcher (Playwright + httpx)
|   |   |-- parser.py        # HTML/content parser
|   |   |-- pdf_extractor.py # PDF text extraction
|   |   |-- frontier.py      # URL frontier management
|   |   |-- robots.py        # robots.txt compliance
|   |   |-- sitemap.py       # Sitemap parser
|   |   |-- url_filter.py    # URL filtering rules
|   |   |-- ingest_client.py # AI-services ingestion client
|   |   |-- tasks.py         # Celery task definitions
|   |   `-- config.py        # Crawler configuration
|   `-- requirements.txt
|
|-- devops/                      # Infrastructure Configuration
|   |-- prometheus/            # Prometheus scrape config
|   `-- grafana/               # Grafana dashboards + provisioning
|
|-- data/                        # Local volumes (gitignored)
|   |-- raw/                   # Raw crawled HTML/PDF
|   |-- processed/             # Processed documents
|   |-- uploads/               # User uploads
|   |-- qdrant/                # Qdrant vector storage
|   `-- redis/                 # Redis persistence
|
|-- docker-compose.yml           # Full stack orchestration (8 services)
|-- Makefile                     # Common dev commands
|-- .env.example                 # Environment variable template
`-- README.md                    # Project documentation
\end{verbatim}

\section*{Key Code Samples}

\subsection*{1. Frontend -- App.jsx (Main Application Component)}

\begin{lstlisting}[language=JavaScript, caption=App.jsx -- React Application with Role-Based Routing, basicstyle=\tiny\ttfamily]
import { Routes, Route, Navigate } from 'react-router-dom'
import { useEffect } from 'react'
import { useThemeStore } from '@/store/themeStore'
import { useAuthStore } from '@/store/authStore'

// Layouts
import PublicLayout     from '@/layouts/PublicLayout'
import AuthLayout       from '@/layouts/AuthLayout'
import MainLayout       from '@/layouts/MainLayout'
import DashboardLayout  from '@/layouts/DashboardLayout'

// Guards
import ProtectedRoute   from '@/components/auth/ProtectedRoute'
import AdminRoute       from '@/components/auth/AdminRoute'
import SuperAdminRoute  from '@/components/auth/SuperAdminRoute'

// Pages
import LandingPage       from '@/pages/LandingPage'
import LoginPage         from '@/pages/LoginPage'
import RegisterPage      from '@/pages/RegisterPage'
import ChatPage          from '@/pages/ChatPage'
import CollegesPage      from '@/pages/CollegesPage'
import CollegeDetailPage from '@/pages/CollegeDetailPage'
import CollegeComparePage from '@/pages/CollegeComparePage'
import AdminDashboard    from '@/pages/admin/AdminDashboard'
import SuperAdminDashboard from '@/pages/superadmin/SuperAdminDashboard'

import { Toaster } from '@/components/ui/toaster'

// Smart redirect based on user role
function RoleRedirect() {
  const { user } = useAuthStore()
  if (user?.role === 'superadmin') return <Navigate to="/superadmin" replace />
  if (user?.role === 'admin')      return <Navigate to="/admin"      replace />
  return <Navigate to="/chat" replace />
}

export default function App() {
  const { theme } = useThemeStore()
  const { initAuth, isInitialising } = useAuthStore()

  useEffect(() => {
    const root = document.documentElement
    root.classList.remove('light', 'dark')
    if (theme === 'system') {
      root.classList.add(
        window.matchMedia('(prefers-color-scheme: dark)').matches ? 'dark' : 'light'
      )
    } else { root.classList.add(theme) }
  }, [theme])

  useEffect(() => { initAuth() }, [initAuth])

  if (isInitialising) {
    return (
      <div className="min-h-screen flex items-center justify-center bg-background">
        <div className="animate-spin rounded-full h-8 w-8 border-b-2 border-primary" />
      </div>
    )
  }

  return (
    <>
      <Routes>
        {/* Public routes */}
        <Route element={<PublicLayout />}>
          <Route path="/"         element={<LandingPage />} />
          <Route path="/colleges" element={<CollegesPage />} />
          <Route path="/colleges/:slug" element={<CollegeDetailPage />} />
          <Route path="/compare"  element={<CollegeComparePage />} />
        </Route>

        {/* Auth routes */}
        <Route element={<AuthLayout />}>
          <Route path="/login"    element={<LoginPage />} />
          <Route path="/register" element={<RegisterPage />} />
        </Route>

        {/* Protected user routes */}
        <Route element={<ProtectedRoute />}>
          <Route path="/dashboard" element={<RoleRedirect />} />
          <Route element={<MainLayout />}>
            <Route path="/chat"            element={<ChatPage />} />
            <Route path="/chat/:sessionId" element={<ChatPage />} />
          </Route>

          {/* Admin routes */}
          <Route element={<AdminRoute />}>
            <Route element={<DashboardLayout />}>
              <Route path="/admin" element={<AdminDashboard />} />
            </Route>
          </Route>

          {/* SuperAdmin routes */}
          <Route element={<SuperAdminRoute />}>
            <Route element={<DashboardLayout />}>
              <Route path="/superadmin" element={<SuperAdminDashboard />} />
            </Route>
          </Route>
        </Route>
      </Routes>
      <Toaster />
    </>
  )
}
\end{lstlisting}

\subsection*{2. AI Services -- LangGraph Agent Nodes (Multi-Agent RAG Pipeline)}

\begin{lstlisting}[language=Python, caption=agents/nodes.py -- Query Understanding and Hybrid Retrieval Nodes, basicstyle=\tiny\ttfamily]
"""Individual MCP-style agent nodes used inside the LangGraph workflow."""
from __future__ import annotations
from typing import Any
from langchain_core.messages import SystemMessage, HumanMessage

from ..core.config import get_settings
from ..rag.llm import chat
from ..rag.prompts import INTENT_SYSTEM, SYSTEM_GROUNDED, USER_TEMPLATE
from ..rag.retriever import retrieve_and_rerank
from ..rag.reranker import rerank
from ..rag.translation import detect_language, atranslate
from .json_utils import extract_json
from .safety import looks_like_injection, sanitize_question
from .state import AgentState


# --- Query Understanding Node ---
async def query_understanding(state: AgentState) -> dict[str, Any]:
    s = get_settings()
    raw = sanitize_question(state.get("user_message", ""))

    if looks_like_injection(raw):
        return {
            "errors": ["prompt_injection_blocked"],
            "intent": "other",
            "needs_retrieval": False,
            "detected_language": "en",
            "working_query": raw,
        }

    # Detect language (supports en, hi, mr)
    requested = state.get("requested_language")
    detected = requested or detect_language(raw)
    if detected not in s.langs:
        detected = "en"

    # Translate to English for retrieval
    working = await atranslate(raw, src=detected, tgt="en") \
              if detected != "en" else raw

    # LLM intent classification
    msgs = [SystemMessage(content=INTENT_SYSTEM), HumanMessage(content=working)]
    try:
        resp = await chat(msgs, temperature=0.0)
        parsed = extract_json(resp) or {}
    except Exception:
        parsed = {}

    return {
        "detected_language": detected,
        "working_query": working,
        "intent": parsed.get("intent", "other"),
        "entities": parsed.get("entities", {}) or {},
        "needs_retrieval": parsed.get("needs_retrieval", True),
    }


# --- Hybrid Retrieval Node ---
async def retrieval(state: AgentState) -> dict[str, Any]:
    if not state.get("needs_retrieval", True):
        return {"reranked": [], "candidates": [], "filters": {}}

    entities = state.get("entities", {}) or {}
    filters = {
        "college": entities.get("college"),
        "state": entities.get("state"),
        "year": entities.get("year"),
    }
    filters = {k: v for k, v in filters.items() if v}

    try:
        reranked = retrieve_and_rerank(
            state["working_query"], filters=filters or None
        )
    except Exception:
        reranked = []

    return {"reranked": reranked, "filters": filters}
\end{lstlisting}

\subsection*{3. AI Services -- Streaming Runner (Parallel Retrieval + Web Search)}

\begin{lstlisting}[language=Python, caption=agents/runner.py -- Streaming Orchestrator with Parallel Execution, basicstyle=\tiny\ttfamily]
"""High-level streaming runner with parallel retrieval."""
from __future__ import annotations
import asyncio
from typing import AsyncIterator
from langchain_core.messages import SystemMessage, HumanMessage

from .nodes import query_understanding, retrieval, _format_context, _top_score
from .state import AgentState
from ..rag.llm import chat, chat_stream
from ..rag.prompts import SYSTEM_GROUNDED, USER_TEMPLATE, HISTORY_TEMPLATE
from ..rag.reranker import rerank
from ..core.config import get_settings


async def run_streaming(
    user_message: str,
    history: list[dict],
    language: str | None,
    web_search: bool | None = None,
) -> AsyncIterator[dict]:
    s = get_settings()
    state: AgentState = {
        "user_message": user_message,
        "history": history or [],
        "requested_language": language,
    }

    # 1) Query understanding
    yield {"type": "status", "step": "understanding"}
    state.update(await query_understanding(state))

    # 2) Resolve coreferences using conversation history
    resolved_query = await _resolve_query(user_message, history)
    state["working_query"] = resolved_query

    # 3) PARALLEL: local retrieval + web search simultaneously
    yield {"type": "status", "step": "searching"}

    async def _local():
        return await retrieval(state)

    async def _web():
        from ..rag.web_search import live_web_chunks
        try:
            return await live_web_chunks(state["working_query"], k=3)
        except Exception:
            return []

    local_task = asyncio.create_task(_local())
    web_task = asyncio.create_task(_web())

    local_result = await local_task
    state.update(local_result)
    local_chunks = state.get("reranked", []) or []
    web_chunks = await web_task

    # 4) Merge + rerank all chunks together
    all_chunks = list(local_chunks) + web_chunks
    if all_chunks:
        merged = rerank(state["working_query"], all_chunks, top_k=s.rerank_top_k)
    else:
        merged = []

    context, citations = _format_context(merged)
    yield {"type": "citations", "citations": citations}

    # 5) Stream generation in target language
    yield {"type": "status", "step": "generating"}
    system_prompt = _system_prompt_for_language(
        state.get("detected_language", "en")
    )
    user_prompt = USER_TEMPLATE.format(
        history_block="", context=context, question=user_message
    )
    msgs = [SystemMessage(content=system_prompt), HumanMessage(content=user_prompt)]

    async for token in chat_stream(msgs, temperature=s.ollama_temperature):
        yield {"type": "token", "text": token}

    confidence = min(1.0, max(0.0, _top_score(merged)))
    yield {"type": "meta", "confidence": confidence, "web_used": len(web_chunks) > 0}
    yield {"type": "done"}
\end{lstlisting}

\subsection*{4. AI Services -- Hybrid Retriever (Dense + BM25 Fusion)}

\begin{lstlisting}[language=Python, caption=rag/retriever.py -- Hybrid Dense+Sparse Retrieval with Reranking, basicstyle=\tiny\ttfamily]
"""Hybrid retrieval (dense + BM25) with metadata filters and reranking."""
from __future__ import annotations
from typing import Optional

from ..core.config import get_settings
from .embeddings import embed_query
from .vectorstore import search as dense_search
from .bm25_index import get_bm25
from .reranker import rerank


def _normalize(scores: list[float]) -> list[float]:
    if not scores:
        return scores
    lo, hi = min(scores), max(scores)
    if hi - lo < 1e-9:
        return [1.0 for _ in scores]
    return [(s - lo) / (hi - lo) for s in scores]


def hybrid_retrieve(
    query: str,
    top_k: int | None = None,
    filters: Optional[dict] = None,
) -> list[dict]:
    s = get_settings()
    k = top_k or s.retrieval_top_k

    # Dense retrieval (Qdrant)
    qvec = embed_query(query)
    dense = dense_search(qvec, k, filters=filters)

    # Sparse retrieval (BM25)
    sparse = get_bm25().search(query, top_k=k, filters=filters)

    # Reciprocal Rank Fusion
    by_id: dict[str, dict] = {}
    dense_scores = _normalize([d.get("score", 0.0) for d in dense])
    for d, ns in zip(dense, dense_scores):
        by_id[d["id"]] = {**d, "_dense": ns}

    sparse_scores = _normalize([d.get("score", 0.0) for d in sparse])
    for d, ns in zip(sparse, sparse_scores):
        existing = by_id.get(d["id"])
        if existing:
            existing["_sparse"] = ns
        else:
            by_id[d["id"]] = {**d, "_sparse": ns}

    # Weighted fusion: alpha * dense + (1-alpha) * sparse
    fused: list[dict] = []
    for c in by_id.values():
        ds = c.get("_dense", 0.0)
        ss = c.get("_sparse", 0.0)
        c["fused_score"] = s.hybrid_alpha * ds + (1 - s.hybrid_alpha) * ss
        fused.append(c)

    fused.sort(key=lambda x: x["fused_score"], reverse=True)
    return fused[:k]


def retrieve_and_rerank(query: str, filters: Optional[dict] = None) -> list[dict]:
    s = get_settings()
    candidates = hybrid_retrieve(query, top_k=s.retrieval_top_k, filters=filters)
    return rerank(query, candidates, top_k=s.rerank_top_k)
\end{lstlisting}

\subsection*{5. Backend -- Conversations API (Streaming Chat with NDJSON)}

\begin{lstlisting}[language=Python, caption=api/conversations.py -- Streaming Chat Endpoint, basicstyle=\tiny\ttfamily]
import json
from fastapi import APIRouter, Depends, HTTPException
from fastapi.responses import StreamingResponse
from motor.motor_asyncio import AsyncIOMotorDatabase

from ..core.database import get_db
from ..core.deps import current_user
from ..core.rate_limit import rate_limit
from ..models.conversation import AskRequest
from ..repositories.conversation_repo import ConversationRepo
from ..services.ai_client import get_ai_client

router = APIRouter(prefix="/api/conversations", tags=["conversations"])


@router.post("/{conv_id}/ask")
async def ask(
    conv_id: str,
    body: AskRequest,
    user: dict = Depends(current_user),
    db: AsyncIOMotorDatabase = Depends(get_db),
    _: None = Depends(rate_limit),
):
    repo = ConversationRepo(db)
    conv = await repo.get_conversation(conv_id, user["_id"])
    if not conv:
        raise HTTPException(404, "conversation not found")

    # Persist user message
    await repo.add_message(
        conversation_id=conv_id, role="user",
        content=body.message, language=body.language
    )

    history = await repo.list_messages(conv_id)
    history_payload = [{"role": m["role"], "content": m["content"]} for m in history]
    ai = get_ai_client()

    if not body.stream:
        result = await ai.ask(
            message=body.message, history=history_payload,
            language=body.language, web_search=body.web_search,
        )
        await repo.add_message(
            conversation_id=conv_id, role="assistant",
            content=result.get("answer", ""),
            citations=result.get("citations", []),
            confidence=result.get("confidence"),
        )
        return result

    # NDJSON streaming response
    async def event_stream():
        accumulated = ""
        citations, confidence = [], None
        try:
            async for chunk in ai.ask_stream(
                message=body.message, history=history_payload,
                language=body.language, web_search=body.web_search,
            ):
                if chunk.get("type") == "token":
                    accumulated += chunk.get("text", "")
                elif chunk.get("type") == "citations":
                    citations = chunk.get("citations", [])
                elif chunk.get("type") == "meta":
                    confidence = chunk.get("confidence")
                yield (json.dumps(chunk) + "\n").encode("utf-8")
        finally:
            if accumulated:
                await repo.add_message(
                    conversation_id=conv_id, role="assistant",
                    content=accumulated, citations=citations,
                    confidence=confidence,
                )

    return StreamingResponse(
        event_stream(), media_type="application/x-ndjson",
        headers={"Cache-Control": "no-cache", "X-Accel-Buffering": "no"},
    )
\end{lstlisting}

\subsection*{6. Docker Compose -- Full Stack Orchestration}

\begin{lstlisting}[language=bash, caption=docker-compose.yml -- Microservices Orchestration, basicstyle=\tiny\ttfamily]
name: arambh

services:
  # --- Infrastructure ---
  redis:
    image: redis:7-alpine
    restart: unless-stopped
    ports: ["6379:6379"]
    healthcheck:
      test: ["CMD", "redis-cli", "ping"]

  qdrant:
    image: qdrant/qdrant:v1.11.3
    restart: unless-stopped
    ports: ["6333:6333", "6334:6334"]
    volumes: ["./data/qdrant:/qdrant/storage"]

  # --- Application Services ---
  backend:
    build: { context: ./backend, dockerfile: Dockerfile }
    depends_on: { redis: { condition: service_healthy } }
    ports: ["8000:8000"]
    command: uvicorn app.main:app --host 0.0.0.0 --port 8000 --reload

  ai-services:
    build: { context: ./ai-services, dockerfile: Dockerfile }
    depends_on: [redis, qdrant]
    extra_hosts: ["host.docker.internal:host-gateway"]
    ports: ["8100:8100"]
    command: uvicorn app.main:app --host 0.0.0.0 --port 8100 --reload

  crawler:
    build: { context: ./crawler, dockerfile: Dockerfile }
    depends_on: [redis, ai-services]
    command: celery -A app.celery_app worker -B -l info --concurrency=2

  frontend:
    build: { context: ./frontend, dockerfile: Dockerfile }
    depends_on: [backend]
    ports: ["5173:5173"]
    command: npm run dev -- --host 0.0.0.0

  # --- Observability ---
  prometheus:
    image: prom/prometheus:latest
    ports: ["9090:9090"]

  grafana:
    image: grafana/grafana:latest
    ports: ["3001:3000"]
    depends_on: [prometheus]
\end{lstlisting}


\subsection*{7. LangGraph Workflow Definition}

\begin{lstlisting}[language=Python, caption=agents/graph.py -- LangGraph StateGraph with Conditional Edges, basicstyle=\tiny\ttfamily]
"""LangGraph multi-agent workflow definition."""
from __future__ import annotations
from langgraph.graph import StateGraph, END

from .state import AgentState
from .nodes import (
    query_understanding, retrieval, live_search,
    generation, validation, followups_node,
)


def build_graph():
    g = StateGraph(AgentState)
    g.add_node("query_understanding", query_understanding)
    g.add_node("retrieval", retrieval)
    g.add_node("live_search", live_search)
    g.add_node("generation", generation)
    g.add_node("validation", validation)
    g.add_node("followups", followups_node)

    g.set_entry_point("query_understanding")

    def _should_retrieve(state: AgentState) -> str:
        return "retrieval" if state.get("needs_retrieval", True) \
               else "generation"

    g.add_conditional_edges("query_understanding", _should_retrieve, {
        "retrieval": "retrieval",
        "generation": "generation",
    })
    g.add_edge("retrieval", "live_search")
    g.add_edge("live_search", "generation")
    g.add_edge("generation", "validation")
    g.add_edge("validation", "followups")
    g.add_edge("followups", END)

    return g.compile()
\end{lstlisting}

\section*{Technology Stack}

\begin{itemize}
\item \textbf{Frontend:} React 18, Vite 5, JavaScript (JSX), Tailwind CSS 3.4, Radix UI, Zustand (state management), React Router 6, i18next (multilingual), Framer Motion, Recharts, Axios
\item \textbf{Backend API:} Python 3.11, FastAPI 0.115, Motor (async MongoDB), JWT Authentication (RBAC -- user/admin/superadmin), Redis (rate limiting + caching), Prometheus metrics
\item \textbf{AI Services:} LangChain 0.3, LangGraph 0.2 (multi-agent orchestration), Ollama \texttt{llama3} (local LLM), HuggingFace \texttt{BAAI/bge-small-en-v1.5} (embeddings), \texttt{cross-encoder/ms-marco-MiniLM-L-6-v2} (reranker)
\item \textbf{Retrieval:} Hybrid Dense (Qdrant) + Sparse (BM25) with weighted fusion, cross-encoder reranking, live DuckDuckGo web search (Perplexity-style)
\item \textbf{Crawler:} Celery + Redis (task queue), Playwright (JS rendering), BeautifulSoup, Trafilatura, pypdf, pdfplumber
\item \textbf{Database:} MongoDB Atlas (document store), Qdrant v1.11 (vector database), Redis 7 (cache + queue)
\item \textbf{Observability:} Prometheus, Grafana, structlog (structured JSON logging), RAGAS evaluation framework
\item \textbf{Orchestration:} Docker Compose (8 services), Makefile automation
\item \textbf{Languages Supported:} English, Hindi, Marathi (direct multilingual generation)
\end{itemize}

\section*{Repository Information}

\textbf{Complete source code is available at:} \\
GitHub: \url{https://github.com/devinxhacker/ArambhAiAdmissionChatbot}

\vspace{5mm}

\textbf{Note:} The complete codebase includes 30+ React pages/components, a full multi-agent RAG pipeline with 6 LangGraph nodes, hybrid retrieval (dense + sparse + web), an async web crawler with Celery workers, JWT-based RBAC authentication, NDJSON streaming responses, multilingual support (en/hi/mr), and comprehensive Docker Compose orchestration. This appendix shows the core architecture and key files for reference.

\vspace{10mm}
\hrule
